\documentclass{article}
\usepackage[T1]{fontenc}
\usepackage{lmodern}
\usepackage{graphicx}
\usepackage{amsmath,amssymb}
\usepackage[a4paper,margin=2cm]{geometry}
\usepackage[absolute,overlay]{textpos}
\setlength{\TPHorizModule}{1mm}
\setlength{\TPVertModule}{1mm}
\usepackage[indonesian]{babel}
\usepackage{hyperref}
\setlength{\emergencystretch}{2em}
% unit: Halaman 54

\begin{document}

% === Halaman 54 (llm) ===

Menurut ad-Durayhim, jenis-jenis \textit{cipher} dapat dikelompokkan ke dalam delapan tipe:
\begin{enumerate}
  \item transposisi,
  \item substitusi,
  \item penambahan atau reduksi jumlah huruf,
  \item penggunaan piranti sandi,
  \item penggantian huruf dengan angka yang diboboti secara desimal,
  \item penyandian huruf dengan menggunakan kata-kata,
  \item penggantian huruf dengan nama generik,
  \item menggunakan simbol atau tanda untuk menyatakan huruf.
\end{enumerate}

\vspace{10mm}

\textcolor{red}{\textit{Cryptology was born among Arabs. They were the first to discover and write down the methods of cryptanalysis.}}
\textcolor{red}{(David Kahn -- Penulis buku: The Code Breaker)}

\end{document}
